\section{Restricted computability and an exact ruler criterion}
\label{sec:restricted}
The onsite instability in Section~\ref{sec:instability} concerns the
particular programmed rulers, not every possible robust exponent.
Several restricted positive results clarify what additional
certificates would suffice. The response observable must be fixed
before making any stability claim.

For a concave ground-energy density $e$ at a transition $g_c$, put
\[
 s_{\rm RT}(t)=e(g_c)+e'_+(g_c)t-e(g_c+t),\qquad
 \rho_{\rm RT}(t)=s'_{\rm RT}(t+),\quad t>0.
\]
After subtracting a vacuum bulk energy, another convention is
$\rho_{\rm vac}(t)=-e'_+(g_c+t)$. They agree when $e'_+(g_c)=0$,
but not at a jump: if $e(g_c+t)=-mt$, then
$\rho_{\rm vac}=m$ has order zero while
$s_{\rm RT}=\rho_{\rm RT}=0$ has no positive logarithmic order.
Moving the critical point does not remove this distinction.
For a convex thermal pressure use
$\rho_{\rm RT}(t)=P'(g_c+t)-P'_+(g_c)$, so that the excess is
relative to the right tangent. These are density and correlation
claims when specified below, not substitutions for the spectral gap.

\subsection{Classical chains and bounded thermal susceptibility}
\begin{proposition}[Finite-range classical chains]\label{prop:chain-cycles}
The zero-temperature energy density of a finite-alphabet
translation-invariant classical chain with fixed finite range and
affine coupling is the minimum of finitely many affine cycle means.
Immediately to the right of every crossing its right-tangent
deficit vanishes. Fixed-width classical strips have the same property.
At finite temperature, bounded finite local energies give an
analytic pressure at every finite coupling.
\end{proposition}
\begin{proof}
For range $R$, form the de Bruijn graph of $V=q^R$ words. An edge
appends one symbol and has its corresponding local energy. Removing
loops from a finite walk leaves a simple path of at most $V-1$
edges. If edge magnitudes are at most $W$ and $m$ is the least
simple-cycle mean, a walk of $n$ edges has energy at least
$nm-2W(V-1)$. Repetition of a minimizing cycle gives the reverse
thermodynamic bound. There are finitely many simple cycles; affine
local weights give affine means. A finite minimum is piecewise affine,
which proves the tangent assertion. Blocking a fixed-width strip
only changes the finite graph. At finite temperature the weighted
de Bruijn matrix is primitive: some common power has every entry
positive. Its Perron root is simple and positive, hence analytic
locally in each finite coupling. Its logarithm gives the pressure.
\end{proof}
This familiar cycle mechanism does not exclude first-order crossings
or every aperiodic minimizing configuration. The range-three chain
paper of 1973 is by Bundaru, Angelescu and Nenciu \cite{isingChains};
Morita's general finite-range ordering work is separate
\cite{moritaChains}. The direct finite-graph proof supplies the
arbitrary-range assertion here. Diverging strip width, hard infinite
energies and zero-temperature transfer singularities are outside the
finite-temperature assertion.

\begin{proposition}[Explicit classical thermal curvature]
\label{prop:thermal-curvature}
Consider bounded finite-range classical interactions on $\Z^d$,
alphabet size $q$, inverse temperature $\beta>0$, and field $-gN$
with $N=\sum_xn(\sigma_x)$ for a nonconstant onsite charge $n$.
On a compact field interval suppose one-site conditional energies
have oscillation at most $B$, and two charge values differ by at
least $\delta_n>0$. If supports have sup-norm diameter at most $R$,
the pressure $P=\lim (\beta|\Lambda|)^{-1}\log Z_\Lambda$ is
strongly convex with constant
\begin{equation}\label{eq:thermal-curvature}
 c_0=\beta(R+1)^{-d}q^{-2}e^{-2\beta B}\delta_n^2>0.
\end{equation}
If on a right interval it is twice differentiable and $P''\leq C$,
then $c_0t\leq\rho_{\rm RT}(t)\leq Ct$ and its exponent is one.
Without the upper bound, any existing finite positive-onset order
lies in $[0,1]$.
\end{proposition}
\begin{proof}
In an open box select one residue class modulo $R+1$. No interaction
meets two selected sites. Given the complement, the selected spins
are independent and every symbol has conditional probability at
least $p_*=e^{-\beta B}/q$. The identity
$\operatorname{Var}(n)=\frac12\sum_{a,b}p_ap_b(n_a-n_b)^2$
gives conditional variance at least $p_*^2\delta_n^2$. Total
variance therefore bounds $\operatorname{Var}(N)$ below by this
number times the selected-site count. The finite pressure has
$P_\Lambda''=\beta\operatorname{Var}(N)/|\Lambda|$.
Subtract the corresponding quadratic and pass its convexity
inequality to the thermodynamic limit. Selected density tends to
$(R+1)^{-d}$, proving \eqref{eq:thermal-curvature} without
interchanging second derivatives and limits.
The two curvature bounds imply
$c_0t\leq P'(g_c+t)-P'_+(g_c)\leq Ct$ by the mean value
theorem and a right limit. Integration gives tangent excess between
$c_0t^2/2$ and $Ct^2/2$. Dividing logarithms gives order one.
With only the lower bound, division by $\log t<0$ gives order at
most one. Right continuity of the bounded convex derivative gives
$\rho_{\rm RT}(t)\to0$, so its finite order is nonnegative.
\end{proof}

Strong convexity and this conditional-variance/sparse-sublattice
method are established by Dobrushin--Nakhapetian
\cite{pressureConvexity}. Equation~\eqref{eq:thermal-curvature}
is an explicit elementary finite-range specialization. If
$0<c\leq c_0$ and $C_1\geq C$ are supplied effective bounds,
\[
 \left|\frac{\log\rho_{\rm RT}(t)}{\log t}-1\right|
 \leq\frac{\max(|\log c|,|\log C_1|)}{|\log t|}
\]
is a computable dyadic modulus on a supplied right interval. It does
not require evaluating an unknown critical point; evaluation in
recentered coordinates would require a name for that point.
A response $t^\alpha$, $0<\alpha<1$, shows why a divergent
susceptibility can evade the exponent-one conclusion.

A conditional upper-bound certificate consists of the identity
$P''=\beta\sum_x\operatorname{Cov}(n_0,n_x)$ and the estimate
$|\operatorname{Cov}(n_0,n_x)|\leq A e^{-m|x|_1}$ uniformly on
the interval. Then
\[
 C=\beta A\left(\frac{1+e^{-m}}{1-e^{-m}}\right)^d.
\]
For a periodic pure phase average the response over a period cell,
using the bound at every origin; mixing translated phases can add
nondecaying covariances. Stable Peierls-phase mixing from contour
theory \cite{contourTheory} requires its own finite-reference-phase
and low-temperature hypotheses. It does not supply uniform
differentiated critical-neighborhood constants for an arbitrary
hierarchical ruler. The observable qualification also matters:
$P(t)=\max(0,t+ct^2/2)$ has total-density jump one but
right-tangent excess $ct$.

\subsection{Algebraic transfer matrices and isolated effective fixed points}
\begin{proposition}[A fixed-dimensional algebraic transfer class]
\label{prop:algebraic-transfer}
Suppose an exact injective matrix product state has fixed bond
dimension, tensor entries with effective algebraic-function branch
descriptions, and an algebraic critical coupling. If its simple
positive leading transfer eigenvalue is $\lambda_0$, its normalized
subleading radius satisfies $0<r(t)<1$, $r(t)\to1$, and the mode
couples to the chosen finite observable, then
$\xi(t)=-1/\log r(t)=t^{-\gamma+o(1)}$ for an effectively
recoverable positive rational $\gamma$.
\end{proposition}
\begin{proof}
Elimination makes all eigenvalue branches algebraic over the
algebraic-number rational function field. Convergent Puiseux
expansions and their effective equality tests
\cite{effectivePuiseux} apply. Conjugation, modulus, division and
finite maxima preserve algebraic branches. A nonzero branch
difference has a fixed leading sign near zero, so the maximizing
branch stabilizes. Thus $1-r(t)=ct^\gamma(1+o(1))$ with
$c>0$ and positive rational $\gamma$, and the first nonzero term
can be computed from the branch data. Finally $-\log r\sim1-r$.
\end{proof}
The overlap hypothesis prevents attributing a measured length to a
dark transfer mode. For an explicit example take
$A_0=I,A_1=tX,A_2=tY,A_3=tZ$ with Pauli matrices. The channel has
eigenvalues $1+3t^2$ on $I$ and $1-t^2$ on each Pauli matrix, so
$\xi\sim1/(4t^2)$. The physical observable
$O=|0\rangle\langle1|+|1\rangle\langle0|$ has connected correlation
\[
 \langle O_0O_k\rangle_c=
 \frac{4t^2}{(1+3t^2)^2}
 \left(\frac{1-t^2}{1+3t^2}\right)^{k-1},\quad k\geq1.
\]
Indeed the inserted channel takes $I$ to $2tX$ and $X$ to $2tI$;
the left fixed functional is $\operatorname{Tr}/2$.
The tensors span $M_2$ for $t>0$. The two-site parent projection
therefore gives a fixed-range parent with the injective uniqueness
property of \cite[Section~4.1.2]{matrixProductStates}. Its entries
are rational functions obtained by the finite Gram inverse. Since
the projection is bounded, the singularity at zero is removable.
This family has zero frustration-free energy throughout, not an
energy singularity or an asserted affine spectral-gap law.
Approximate MPS, growing bond dimension and unrestricted computable
parameter dependence are outside the proposition.

\begin{proposition}[An isolated RG certificate]\label{prop:rg-certificate}
Let $R$ and $DR$ be computable near a specified rational compact
finite-dimensional box $K$, and suppose $R(x)-x$ has exactly one
zero $x_*$ in $K$. Then $x_*$ and $DR(x_*)$ are computable. A
specified isolated real eigenvalue $\lambda>1$ and a physical
identity $\nu=\log b/\log\lambda$, integer $b\geq2$, compute $\nu$.
\end{proposition}
\begin{proof}
The norm $\|R(x)-x\|$ is computable. Refine rational boxes and
discard a box only after certifying strict positivity throughout it.
Never discard the zero. Outside any neighborhood of the unique zero,
compactness gives a positive minimum, eventually certified by
effective continuity. Dovetail subdivision until the unexcluded
boxes have an enclosing box of the requested small diameter. This
computes $x_*$, including a boundary zero, and then $DR(x_*)$.
Equivalently apply \cite[Theorem~2.9]{invariantComputability}.
Computable polynomial-root approximation gives the spectrum as a
multiset, and the supplied isolator selects $\lambda$. Positivity
of $\log\lambda$ permits the stated division.
\end{proof}
An ordinary $C^1$ promise does not supply an effective derivative.
Neither stability alone nor an arbitrary finite RG truncation
supplies this exact reduction, compact isolator or physical identity.
Computable RG steps can have uncomputable flows
\cite{renormalisation} without contradicting this isolated readout.

\subsection{Why soft thermodynamic data do not suffice}
\begin{proposition}[A volume-grid imitation]\label{prop:polyhedral-imitation}
For each limit-computable $\zeta\in[1,2]$ there is a computable
nondecreasing density $\rho$, zero on $(-\infty,0]$ and one on
$[1,\infty)$, of logarithmic order $\zeta$ at zero. Put
$s(u)=\int_0^u\rho(v)\,dv$ and
$E(A,B;g)=-A s((g-B)/A)$, $A>0$.
For $N=L^d$, uniformly computable rational polyhedral energies $E_L$
with $N+1$ affine branches have volume-grid slopes and satisfy
\[
 E-A/N\leq E_L\leq E+(2A+|g-B|)/N.
\]
They are concave and positively homogeneous in $(A,B,g)$, and their
limiting density order is the constant $\zeta$ on the open domain
$A>0$, $B\in\mathbb R$.
\end{proposition}
\begin{proof}
Take computable rationals $\zeta_j\to\zeta$ in $[1,2]$.
At $t_n=2^{-n}$ prescribe
$\rho(t_n)=2^{-Q_n}$, $Q_n=\sum_{j\leq n}\zeta_j$, with
$Q_0=0$, and interpolate linearly. The ratios of successive values
are in $[1/4,1/2]$, $\rho(t)\leq t$, and the interpolation has a
known bounded Lipschitz constant. The tail integral below $2^{-n}$
is at most $2^{-2n}/2$, so $s$ is computable without a modulus for
$\zeta_j$. Ces\`aro convergence gives $Q_n/n\to\zeta$; bounded
ratios extend the order from nodes to all $t$. Differentiation of
the perspective gives the claimed density order at $g_c=B$.

On $0\leq r\leq1$ let $c(r)=\max_{0\leq u\leq1}(ru-s(u))$.
The compact computable optimization has a known Lipschitz bound,
so $c$ is computable, $0\leq c\leq1$, and
$|c(r)-c(r')|\leq|r-r'|$. Convex duality, or the tangent of slope
$\rho(u)$ with the endpoint slopes on the linear extensions, gives
$-s(u)=\min_{0\leq r\leq1}(c(r)-ur)$ for all real $u$.
For $r_j=j/N$, approximate $c(r_j)$ within $1/(4N)$, round to the
nearest $1/N$ grid and clamp to $[0,1]$, obtaining rational $a_j$
with error at most $1/N$. Set
\[
 E_L=\min_{0\leq j\leq N}\{Aa_j+(B-g)j/N\}.
\]
Rounding lowers no branch by more than $A/N$. Moving a dual
minimizer to the grid costs at most $(A+|g-B|)/N$, and rounding
then costs at most another $A/N$. This proves both bounds.
\end{proof}
Thus computable finite energies, rational volume-grid slopes,
concavity, homogeneity, a boundary-to-volume rate and constancy of
an exponent on an open parameter domain do not force a computable
exponent. The missing microscopic requirement is one common fixed
finite local table whose configurations obey consistent restriction
and gluing rules. The independently rounded branches do not supply
that table. Thermal leading onsite responses with $\zeta>1$ would
also contradict Proposition~\ref{prop:thermal-curvature}; adding an
analytic response changes the leading observable.

One simple sufficient transfer principle is still useful. If an
exponent is constant on an open set and one member has a valid
finite certificate for a computable value, equality transfers that
value to every member. Uniform computation needs the certificate or
an effective guaranteed search. Density of rational interaction
tables does not give density of certified computable exponents.
A more explicit certificate consists of positive computable $X(t)$,
effective $c,C,t_0>0$ and $m(k)$ such that
\[
 ct^{\alpha+2^{-k}}\leq X(t)\leq Ct^{\alpha-2^{-k}}
 \quad(0<t\leq\min(t_0,2^{-m(k)})).
\]
Then the normalized logarithm differs from $\alpha$ by at most
$2^{-k}+\max(|\log c|,|\log C|)/|\log t|$. Choose a small dyadic
$t$ effectively and evaluate its logarithm, finding a positive lower
enclosure first. This computes $\alpha$. It assumes neither a
computable derivative of every computable pressure nor that arbitrary
continuous critical recentering is effective.

\subsection{An exact calibrated ruler and its perturbations}
\label{subsec:exact-ruler}
Fix a finite-dimensional local-potential basis and subtract its
vacuum bulk and linear wall contributions. Assume that the
\emph{full} thermodynamic energy, for every perturbation in the
neighborhood in question, is exactly
\begin{equation}\label{eq:exact-ruler}
 e_v(g)=\min\left(0,\inf_{p\in\mathcal P}
 \left[-\frac gp+\frac{C(p)+v\cdot R(p)}{p^2}\right]\right).
\end{equation}
Here $C(p)>0$, $\log C(p)/\log p\to a\in(0,1)$,
$\|R(p)\|_1=o(p)$, and the unbounded pitch set is logarithmically
dense: the least available $p\geq x$ has $\log p/\log x\to1$.
In a pattern-count interpretation the complete vector is
$\alpha p^2+\beta p+R(p)$; its bulk term has been subtracted and
the physical field is $g+v\cdot\beta$. Equation~\eqref{eq:exact-ruler}
is an assumption about all competing configurations, not just clean
cells. Establishing it for one finite local rule is a separate obligation.

\begin{lemma}[Logarithmically dense ruler optimization]
\label{lem:ruler-optimization}
If $D(p)>0$ has logarithmic order $a<1$ on this pitch set, then
\[
 s(t)=\sup_p\{t/p-D(p)/p^2\}=t^{q+o(1)},\quad
 q=1+\frac1{1-a},\qquad
 s'(t+)=t^{1/(1-a)+o(1)}.
\]
\end{lemma}
\begin{proof}
For small $\delta>0$ and all sufficiently large $p$,
$p^{a-\delta}\leq D(p)\leq p^{a+\delta}$. Each fixed finite pitch
has negative profit for small $t$. Maximizing
$t/p-p^{a-\delta-2}$ over all positive real $p$ bounds $s$ above
by a constant times $t^{1+1/(1-a+\delta)}$.
For a lower bound choose $x=(2/t)^{1/(1-a-\delta)}$ and the least
available $p\geq x$. Then $D(p)/p\leq t/2$ and
$\log p/\log x\to1$, so $s(t)\geq t/(2p)$ has order
$1+1/(1-a-\delta)$. Squeeze logarithmic orders and let
$\delta\downarrow0$. Convexity and $s(0)=0$ give
$s(t)/t\leq s'(t+)\leq s(2t)/t$, proving the derivative order.
The supremum is nonnegative because the branches tend to zero.
\end{proof}

\begin{theorem}[Residual-to-cost criterion]\label{thm:ruler-stability}
Within \eqref{eq:exact-ruler}, a sufficiently small full ball of
perturbations preserves the positive density-onset exponent
$1/(1-a)$ at the actual calibrated transition if and only if
\begin{equation}\label{eq:ruler-ratio}
 \sup_{p\in\mathcal P}\frac{\|R(p)\|_1}{C(p)}<\infty.
\end{equation}
If the ratio is unbounded, every perturbation ball contains one
fixed signed-coordinate perturbation giving a finite-pitch
first-order onset. Its vacuum-relative order is zero and its
right-tangent positive order is undefined.
\end{theorem}
\begin{proof}
If the bound is $K$, choose $\|v\|_\infty<\eta$ with
$\eta K<1/2$. Then $D_v=C+v\cdot R$ lies between $C/2$ and
$3C/2$. It has order $a$, is positive and obeys $D_v/p\to0$.
The vacuum wins for $g\leq0$, and every $g>0$ beats it at large
enough pitch. Thus $g_c=0$, and
Lemma~\ref{lem:ruler-optimization} gives the common exponent.
The right derivative at zero is zero, so both response conventions
agree. If $K=0$ the conclusion is immediate.

If the ratio is unbounded, finite dimensionality supplies a fixed
coordinate and sign $\sigma$ with $\sigma R_i/C$ unbounded above
along a subsequence. Choose any sufficiently small fixed rational
$\varepsilon>0$ and $v=-\varepsilon\sigma e_i$.
Then $D_v$ is negative somewhere. The thresholds
$h(p)=D_v(p)/p$ tend to zero, so their negative infimum $h_*$ is
attained among finitely many pitches. Let $p_*$ be the smallest
minimizer. For $g=h_*+t$ and $p\geq p_*$,
\[
 \frac{h(p)-h_*-t}{p}\geq-t/p\geq-t/p_*.
\]
For the finitely many $p<p_*$ the same comparison holds whenever
$t\leq(h(p)-h_*)/(1-p/p_*)$. A positive common right interval
therefore exists. Below $h_*$ the vacuum wins, while on that
interval $e_v(h_*+t)=-t/p_*$. This is the actual shifted onset,
with jump $1/p_*$. Subtracting its right tangent leaves zero.
\end{proof}
Pattern-discrepancy control relative to defect costs has an important
antecedent in Mi\k{e}kisz's strict-boundary stability analysis
\cite{stableQuasicrystals}. The recentered finite-pitch conclusion
above additionally uses the exact envelope and its tail behavior.
It is not a proof that a full robust lattice realization exists.

An integer-consistent example makes the failure explicit. Let
$p_n=n^4$, $C(p_n)=n$, and formal pattern counts
\[
 (n^8-n^2-n,n^2,n),\qquad R(p_n)=(-n^2-n,n^2,n),\quad n\geq2.
\]
They are nonnegative and sum to $p_n^2$, while
$\|R\|_1/C=2n+2$. The unperturbed density order is $4/3$.
Biasing the second pattern by $-\varepsilon$, with
$\varepsilon=3/(2N)$ and integer $N\geq3$, gives
$h(n)=n^{-3}-\varepsilon n^{-2}$ and the exact identity
\[
 h(n)-h(N)=\frac{(n-N)^2(n+2N)}{2N^3n^3}\geq0.
\]
Hence $g_c=-1/(2N^3)$, $p_*=N^4$ and the density jump is $N^{-4}$.
The perturbation norm tends to zero as $N\to\infty$, but each fixed
perturbation preempts the positive onset. These counts do not
constitute a common finite local rule.

\subsection{Static obstructions to shortcuts from noisy computation}
The following exact checks separate the reliable trajectory theorem
from an equilibrium energy estimate.
First, on a trajectory cloned across a transverse $W\times W$
torus, every full-symbol majority agrees with the base symbol.
The stacked prediction is the base prediction, even for erroneous
base histories. Therefore both update-error counts and any stipulated
replicated additive marks are multiplied by exactly $W^2$, including
differences between two histories. Transverse cloning alone does
not improve their energy ratio.

Second, start a majority history with an all-one plane. Change the
entire next plane to zero and retain zero thereafter. This costs
exactly $JW^2$ in update errors. A positive mark cost $r$ on ones
charges the honest history $rW^2T$, so the competitor wins for
$T>J/r$. There is no uniform estimate bounding all time-distributed
marks by a constant times these faults. This organized deterministic
plane error does not contradict stochastic low-noise reliability.

Third, the blank/nonblank trajectory model of
\cite[equation~(13), Sections~3.2 and~4]{thermalTrajectories} has,
after a constant subtraction, update cost $JD$ with
$J=\log(16(1-\varepsilon_0)/\varepsilon_0)>0$. It has an all-blank
zero-error trajectory and fully nonblank zero-error trajectories
from valid tilings. If $Q$ counts nonblank sites, adding blank and
nonblank fields and subtracting the vacuum constant gives
$H=JD+(\varepsilon-g)Q$. Since $D\geq0$ and
$0\leq Q\leq|\Lambda|$, its energy density is bounded below by
$\min(0,\varepsilon-g)$. The two zero-error endpoints attain the
bound; restricting to open boxes costs only boundary terms. Thus
\[
 e(g)=\min(0,\varepsilon-g).
\]
This is a density jump of one, not the positive tangent onset sought
above. At $g=0$, every translation-invariant ground-state measure
has zero expected $JD+\varepsilon Q$ and hence is all blank.
This measure assertion is not uniqueness of every configuration
under finite modifications.

Finally the desired arithmetic can be scheduled without thereby
solving the static construction. For computable rationals $a_k\to
a\in(0,1)$, set $s_0=0$,
$d_k=[s_{k-1}<ka_k]$, $s_k=s_{k-1}+d_k$,
$p_k=N^k$, $m_k=N^{s_k}$ for fixed $N\geq2$.
For any eventual band $0<L<a_k<U<1$, induction gives
$Lk-C\leq s_k\leq Uk+C$: a zero step certifies the lower bound,
and a one step certifies the upper bound; the opposite bounds
propagate. Hence $s_k/k\to a$ and $\log m_k/\log p_k\to a$.
Choosing $N$ children or one child at stage $k$ gives this count.
For a slow rational generator, simulate it for $k$ steps and retain
the latest completed rational. Completed indices still tend to
infinity and finite-string feedback comparisons have polynomial
overhead in $k$. This supplies a counting schedule, not protected
calibrated minimum costs, an interaction compiler or the full
energy identity \eqref{eq:exact-ruler}.

These results complement the explicit instability of the programmed
rulers and the instability theorem of Castilla-Castellano--Lucia
\cite{gapInstability}. Neither instability result proves a general
equivalence between robustness and computability. The restricted
certificates above identify sufficient mechanisms and the exact
ruler theorem identifies its own necessary cost control; the general
equilibrium question remains separate.

\subsection{Anchored robust density and counting laws}
\label{subsec:robust-laws}
The reliable computation of Section~\ref{subsec:reliable-readout}
also supports a fixed scale readout with an exact arithmetic range.
The anchors are essential: a computable bottom configuration, a
distinguished program origin, infinite cloning in two directions,
and a prescribed time-then-window limit. This is a theorem about
anchored causal histories, not an intrinsic equilibrium gap exponent.

Fix a finite alphabet $\Sigma$, a finite neighborhood and a
positive product-update kernel $K$. Its simplex has one probability
vector for each predecessor pattern. A law is robust on a nonempty
relatively open set $U$ when, for each fixed initial configuration
under consideration, it exists and is constant as $K$ varies in $U$.
Different programs may have different robust laws. At computable
kernel entries and a computable initial configuration, every finite
spacetime cylinder probability is computable by its finite ancestor
graph. All supports, site maps and cylinder functions below are
uniformly computable from their integer indices: one algorithm returns
each finite support and computes its cylinder function to any requested
rational accuracy. This includes the canonical geometry and projection,
which are fixed independently of the program and kernel. A fixed but
noncomputable choice of sites would instead supply extra advice.

At even $t\geq2$ choose the odd, computable schedule
\[
 m_t=2\lceil\log_2t\rceil^2+1.
\]
The ceiling here is computed by integer bit length. This rounded
logarithmic-square schedule has the same summability and
growth-versus-mixing properties used below.
For a fixed bit projection $b$ and computable site map $z(n,t)$,
let $O(n,t)$ be the majority of $b(X_t)$ at
$z(n,t)+\ell(2t+1)e_1$, $0\leq\ell<m_t$.
The redundant-coordinate speed of the stack is at most one, so these
endpoints have disjoint depth-$t$ noise cones. Put
\begin{align}
 \rho_N(t)&=N^{-1}\sum_{n<N}O(n,t),&
 c_N(t)&=\sum_{n<N}O(n,t),\notag\\
 \Phi_{\rm dens}&=\lim_{N\to\infty}\lim_{t\to\infty}
                 \mathbb E\rho_N(t),&
 \Phi_{\rm exp}&=\lim_{N\to\infty}
  \frac{\log(1+\lim_{t\to\infty}\mathbb E c_N(t))}{\log N}.
 \label{eq:robust-observables}
\end{align}
Use $N\geq2$ in the exponent, and restrict to existing indicated
limits. In particular $1+$count, rather than the count itself, is
strictly positive. The sample analogues use the same prescribed
order. They are not assertions that one trajectory executes the
outer limit in time. A depth-$d$ readout means $d$ specified nested
existing limits of expectations of uniformly computable finite
cylinder functions, with any stated computable continuous operations.
This broad decoder class can itself carry a target program: a constant
cylinder function $N^{-1}\sum_{n<N}a(n,t)$ can ignore the world and
have any $\mathsf D_3$ density as its two-limit value when the decoder
varies with the target. The result below instead fixes one counting
readout before choosing the finite program in the dynamics. It does
not classify every computable decoder by a memory mechanism.

\begin{proposition}[Robust depth ceiling]\label{prop:robust-depth}
At computable data and kernel, a depth-$d$ value is in
$\mathsf D_{d+1}$. If it is robust on a nonempty open kernel set,
the same absolute ceiling holds throughout that set, including its
noncomputable kernels.
\end{proposition}
\begin{proof}
Replace each finite expectation by rational approximations with
effective vanishing errors. One existing limit is computable in
$0'$ by Cauchy-tail search; iterate relative to successive jumps.
For \eqref{eq:robust-observables}, the nonnegative inner count has
a $0'$-name, which computes the logarithm of one plus it. A further
existing limit gives $0''$. Every nonempty relatively open region
of the finite positive-rate simplex contains a computable kernel.
A robust value equals its value there; this is witness transfer,
without a uniform algorithm locating a witness in an arbitrary open
set presentation.
\end{proof}

\begin{theorem}[Anchored robust third level]\label{thm:robust-third}
One fixed computable universal three-dimensional stacked rule $G$,
one nonempty open positive-rate kernel region $U$, and one fixed
computable projection and site map in \eqref{eq:robust-observables}
have the following property. Starting from the computable bottom
configuration obtained by infinitely cloning a finite program at
a distinguished origin, the ranges over programs of the existing
density laws and of the existing counting-exponent laws are each
exactly $\mathsf D_3\cap[0,1]$, for every $K\in U$.
Each law is independent of $K$, and its expectation and sample
iterated limits agree almost surely. All these limits retain the
prescribed time-then-window order and the stated anchors.
\end{theorem}
\begin{proof}
Fix a universal ordinary interpreter with a product tape alphabet:
a protected output-bit track, scratch tracks and head/state fields.
Every ordinary instruction copies the output track. Only a designated
write instruction changes a selected output bit, including when a
head passes through its cell. Initialize all output bits to zero,
the finite program and head near the origin, and other scratch
cells blank. A fixed radius-$r$ local CA simulation preserves this
track invariant. At stage $s$ the interpreter computes a total
binary approximation $a(0,s),\ldots,a(s,s)$ and writes register $n$
at original coordinate $-2r(n+1)$. If $a(n,s)\to A(n)$, infinitely
many completed stages visit $n$ and all sufficiently late writes
repeat $A(n)$. Between writes the output stays fixed. Thus its
logical bit $L_n(t)$ stabilizes, without an effective settling bound.

For shift $(Sx)_i=x_{i+1}$, set $H=S^rF$ and group $2r$ cells.
Then $H^t(x)_i=F^t(x)_{i+rt}$ and the blocked neighborhood is
$\{0,1\}$. At even times the register is at blocked coordinate
$-(n+1)-t/2$, internal offset zero. Hence
\[
 z(n,t)=(0,0,-(n+1)-t/2)
\]
and one fixed bit projection work for every program. Stack this
blocked rule using full-symbol north-east-center majority. The
Toom/G\'acs--Reif reliability theorem, in the arbitrary-trajectory
form \cite[Theorem~3.1]{thermalTrajectories}, gives a positive
rational $\varepsilon_0$ for site error target $1/4$. Define
\[
 U=\{K: K(a\mid r)>0\text{ for all }a,r,\quad
              K(G(r)\mid r)>1-\varepsilon_0\text{ for all }r\}.
\]
The rule, region and readout are fixed before choosing the program.
Product updates give the required finite joint-fault bounds by
conditioning backwards over time layers. The cloned bottom is
computable, even though its redundancy is infinite.

The $m_t$ endpoint symbols use disjoint independent noise variables.
A bit error requires a symbol error, of probability at most $1/4$.
Hoeffding therefore gives
\begin{equation}\label{eq:robust-amplification}
 \Pr\{O(n,t)\ne L_n(t)\}\leq e^{-m_t/8}=:e_t,
 \qquad |\mathbb E O(n,t)-L_n(t)|\leq e_t.
\end{equation}
The bound is summable over even $t$, since $m_t$ grows as
$(\log t)^2$. Borel--Cantelli gives eventual equality to the current
logical bit for each fixed $n$. A countable intersection covers all
registers. Hence the inner limits of every fixed window, both in
expectation and almost surely, are exactly its count in $A$.
The probability-one event can depend on $K$.

For any $\beta\in\mathsf D_3\cap[0,1]$, take uniformly
$0'$-computable rationals $u_s\to\beta$ in $[0,1]$. With
$N_0=1$, $A(0)=0$, $N_{s+1}=(s+2)N_s$, put the mechanical bits
$A(N_s+j)=\lfloor(j+1)u_s\rfloor-\lfloor ju_s\rfloor$ on the
successive blocks. If $d_s=|A\cap[0,N_s)|/N_s$, then for
$N=N_s+j$,
\begin{equation}\label{eq:robust-all-prefix}
 \left|\frac{|A\cap[0,N)|}{N}-\beta\right|
 \leq\frac{N_s}{N}|d_s-\beta|
      +\frac jN|u_s-\beta|+\frac1N.
\end{equation}
At the next endpoint the first weight is $1/(s+2)$, so endpoint
errors tend to zero. Equation~\eqref{eq:robust-all-prefix} then
proves convergence at every intermediate prefix, including
$\beta=0,1$. The set is $0'$-computable. The ordinary limit lemma
supplies the total computable approximation executed by the interpreter.
This density correspondence is classical
\cite{density,ershovDensity}; it is not a new density theorem.

For the counting exponent $\gamma\in\mathsf D_3\cap[0,1]$ use
$0'$-computable rational $u_k\to\gamma$ and
\[
 A=\bigcup_{k\geq1}[2^k,2^k+M_k),\qquad
 M_k=\lfloor2^{ku_k}\rfloor.
\]
For $u_k=a/b$, $M_k$ is the largest integer $m$ with
$m^b\leq2^{ka}$, computed by exact integer comparison. Thus a
floor at integer equality does not require a new real-equality
oracle. If $2^K\leq N<2^{K+1}$, the full preceding block gives
lower bound $M_{K-1}$ and the sum through block $K$ gives an upper
bound. Eventual $|u_k-\gamma|<\eta$ yields
\[
 (K-1)(\gamma-\eta)-O(1)
 \leq\log_2(1+|A\cap[0,N)|)
 \leq K(\gamma+\eta)+O(\log K)+O(1),
\]
with zero as the lower bound if $\gamma<\eta$. Divide by
$\log_2N\in[K,K+1)$ and let $\eta\downarrow0$. This includes
both endpoints. Compile this set's approximation in the same way.
Conversely, suppose all inner prefix expectation limits exist at any
$K\in U$, including a noncomputable kernel. The amplification error
in~\eqref{eq:robust-amplification} tends to zero, so each deterministic
integer prefix count $\sum_{n<N}L_n(t)$ converges along the sampled
even times and is eventually constant. Differences of consecutive
prefix counts make the corresponding logical output bits eventually
constant along those times. If the exponent starts at $N=2$, choose
the first two bits to have that prefix's eventual count; this finite
choice changes neither asymptotic statistic. The resulting set is
limit-computable, independently of $K$. Every stipulated fixed-window
expectation and sample limit is its count, uniformly throughout $U$.
Its existing density or counting order is in
$\mathsf D_3$ by the same relativized limit argument. Thus necessity
does not assume robustness for an arbitrary program at a noncomputable
kernel. Proposition~\ref{prop:robust-depth} supplies the corresponding
ceiling at a computable witness.
\end{proof}
The programs realizing a prescribed density and a prescribed
counting exponent may differ. A positive-density set has counting
exponent one; the theorem does not assert arbitrary independent
pairs for one program. The fixed-rule synthesis uses the prior
reliable-computation technology \cite{toom,gacsReif,thermalTrajectories}
and the asymptotic-density literature \cite{density,ershovDensity}.
At fixed finite time, only finitely many logical output registers can
have been reached from the program. The remaining far-blank region
has one common noisy-majority expectation $q_t(K)\leq e_t$, rather
than an exactly zero physical value. Thus
$\lim_t\lim_N\mathbb E\rho_N(t)=0$, whereas the prescribed
opposite order is $\beta$. For $\beta\ne0$ the two orders differ.

\subsection{Prescribed higher depths and restricted memory mechanisms}
\begin{corollary}[Anchored higher-depth readouts]\label{cor:robust-depth}
For each fixed $d\geq2$, the same protected interpreter, kernel
region and anchored initialization realize exactly
$\mathsf D_{d+1}\cap[0,1]$ by a fixed depth-$d$ readout with
$d-2$ tuple-index limits outside time and window. A separate fixed
one-limit weighted readout realizes $\mathsf D_2\cap[0,1]$.
The computable bottom boundary, distinguished program origin,
infinite cloning and prescribed outer limits are retained.
\end{corollary}
\begin{proof}
Iterated relative limit approximation expresses a target as
$d-2$ existing limits of uniformly $0''$-computable real names
$c(k_1,\ldots,k_{d-2})\in[0,1]$. The conversions must use names,
not uniformly decidable strict rational cuts. Explicitly, for
$0''$-computable rational approximants $q(k,j)$ with error $2^{-j}$,
the relative limit lemma gives $0'$-computable guesses
$q_s(k,j)$ eventually equal to each rational code. At stage $s$
choose the largest $m\leq s$ for which the first $m$ closed rational
error intervals meet $[0,1]$, and output a rational in that intersection;
use a default if necessary. Each fixed finite prefix is eventually
correct, so $m\to\infty$ and the outputs converge to $c(k)$.
This avoids an unjustified diagonal approximation to the jump.
Apply the mechanical blocks uniformly in the tuple, then the
uniform ordinary limit lemma to obtain $a(k,n,s)$. Fix a computable
encoding of tuples into register indices and service the first $s$
encoded registers at stage $s$. Every register is visited infinitely
often. The prescribed outer limits now give the target. Conversely,
at any kernel where the stipulated limits exist, the consecutive-prefix
argument above applies uniformly to every tuple. It gives a
limit-computable logical table whose inner values agree throughout
$U$; the depth ceiling then gives necessity. The readout is fixed separately for
each $d$; changing depth changes the definition, not the physical rule.
For $d=1$, use $\sum_{n<t}2^{-n-1}O(n,t)$. A $0'$-computable
binary expansion of a target has a computable stabilizing bit
approximation. At dyadic targets choose a fixed computable expansion;
for target one use the all-ones fractional expansion.
Finite-prefix stabilization and the geometric tail
give expectation and almost-sure convergence. This is a one-limit
readout, not the two-limit density renamed.
\end{proof}

Three restricted comparisons are useful, without being exhaustive.
For a kernel with Dobrushin coefficient $\kappa=\sum_vc_v<1$,
sitewise coupling gives marginal distance at most $\kappa^t$ from
its unique stationary law. Let its projected bit marginal be $\pi$.
Disjoint-cone endpoint samples from deterministic data remain
independent. If $\pi\ne1/2$, an eventual fixed bias and majority
amplification give expectation limit zero or one. If $\pi=1/2$,
couple the $m_t$ bits to independent fair bits to obtain error at
most $m_t\kappa^t\to0$. The density expectation law is therefore
in $\{0,1/2,1\}$ and its counting exponent in $\{0,1\}$.
At a fixed kernel the initial-data image is a singleton; the three
possibilities are across kernels. No almost-sure convergence of a
Boolean readout to $1/2$ is asserted. For a general mixing error
$\delta_t$, the stronger condition $m_t\delta_t\to0$ suffices;
the sharper condition $\sqrt{m_t}\delta_t\to0$ also follows from
binomial comparison. Arbitrarily growing majorities need not forget.

For static Toom memory, use the stationary site map
$z_{\rm id}(n,t)=(0,0,-(n+1))$ for an unshifted binary identity
rule, with the same transverse copy spacing. Alternatively shift and
block both the identity dynamics and the readout together. Using
the moving universal-computation site map with an unshifted identity
would sample different tape bits. With the matching coordinates,
every computable tape is recovered with error $e_t$, so its existing
density lies in $\mathsf D_2$; computable mechanical tapes attain
every such value. In the same universal rule as
Theorem~\ref{thm:robust-third}, programs with a computable logical
settling bound $S(n)$ likewise give exactly $\mathsf D_2$ densities:
$A(n)=L_n(S(n))$ is computable, and write-once total programs realize
all computable tapes. Here $S(n)$ is a bound in physical CA time;
a bound given in interpreter stages is converted by deterministic
simulation until the corresponding commits finish. Choose computable
even $T(N)$ beyond the first
$N$ settling bounds with $Ne_{T(N)}\leq2^{-N}$. The observed mean
then converges to its density almost surely by a union bound and
Borel--Cantelli. This is a data-reading realization, with the anchors
and growing support retained.

At a computable kernel and computable data, a supplied computable
inner modulus $S(n,j)$ makes
the inner expectations uniformly computable; one remaining existing
limit is in $\mathsf D_2$. Such effectiveness is pointwise unless
uniform modulus data are supplied. A name of a computable kernel
does not uniformly decide whether $\pi=1/2$. These are sufficient
mechanisms, not an exhaustive classification or a necessary condition
for a particular law to be computable. For example, put halting-set
guesses only at register indices $2^j$, all other bits zero. The
eventual set has density zero but no computable uniform settling
bound. Robust equality transfers the ceiling from a computable kernel
to the whole region. The Ces\`aro and full-Gibbs-boundary moduli of
Proposition~\ref{prop:unique-moduli} give additional restricted
effective-inner-limit mechanisms.
For equilibrium responses, Section~\ref{sec:ceilings} computes pressure
and first derivatives at computable points of promised differentiability.
With a computable critical point and tangent, a strictly positive
computable tangent deficit has an existing logarithmic order in
$\mathsf D_2$. A positive response $f(t)=t^{a+o(1)}$, $a>-1$,
has integral of order $a+1$: bound it by $t^{a\pm\eta}$ and
integrate, taking $\eta<a+1$. This transfers that conditional
deficit-order ceiling to the corresponding response order. The
critical data and positivity are substantive hypotheses. The known
unique-measure mechanism is credited in Theorem~\ref{thm:unique-states}.

\subsection{Rates and robust hidden settling scales}
\begin{proposition}[Rates with computable reference scales]
\label{prop:robust-rates}
Let uniformly computable $q_s\to r$. If $r$ is noncomputable, every
positive computable $h(s)\to0$ is exceeded by $|q_s-r|$ infinitely
often. Against a positive computable reference $h_0(s)\to0$, any
existing finite exponent $\lim_s\log|q_s-r|/\log h_0(s)$ is zero.
If $r$ is computable and the errors are eventually nonzero, that
existing finite exponent is in $\mathsf D_2$.
\end{proposition}
\begin{proof}
An eventual error bound $h(s)$, with a finite cutoff hard-coded,
would compute $r$ by searching for sufficiently small $h(s)$ and
approximating $q_s$. A positive rate exponent would supply such
a bound $h_0(s)^a$ for a positive rational $a$. Since the error
tends to zero, the defined quotient is eventually nonnegative.
For computable $r$, positive errors and their logarithms are
computable on the promised domain, so the exponent is an existing
limit of computable values. Zeros need a separate convention and
are not passed silently through a real logarithm.
\end{proof}

Fix a program in Theorem~\ref{thm:robust-third} with a density or
counting exponent $\beta\notin\mathsf D_2$. Let $\tau(N)$ be the
least time after which its first $N$ logical bits never change.
This deterministic quantity is identical for every $K\in U$.
For every total computable $B$, infinitely many $N$ have
$\tau(N)>B(N)$: an eventual bound, with the finite cutoff hard-coded,
would compute every final bit and make the existing law
$\mathsf D_2$. Consequently the usual power, exponential and
fixed-height iterated-exponential growth comparisons have no finite
upper exponent; an extended value $+\infty$ is not excluded.

After $t\geq\tau(N)$, expectation error of the count is at most
$Ne_t$, and the probability of any erroneous bit in the window is
at most $Ne_t$. A sufficient accuracy time is the maximum of
$\tau(N)$ and a computable amplification margin. This is not an
equality of settling times or a deterministic last-fault bound.
At each $K\in U$, even a physical guarantee of window-density error
below $1/(4N)$ for every $t\geq B(N)$ cannot have a computable
eventual upper bound. Otherwise choose computable even $t\geq B(N)$
with $e_t<1/(4N)$. The current logical prefix density is within
$1/(2N)$ of the final one. Both lie on the $1/N$ grid, so their
counts are equal. This computes all final prefix counts and hence
the noncomputable set.

At \emph{computable} $K$, any computable even diagonal $\sigma(N)$
gives computable finite expectations. An existing diagonal limit
is $\mathsf D_2$, hence cannot be $\beta$. The corresponding sample
diagonal, using rational approximations for the logarithmic statistic
when necessary, is an effective inquiry; Theorem~\ref{thm:reach}
makes its fixed-target convergence probability zero. The physical
law itself is still the specified iterated limit.

For the density there is also a pointwise statement at every $K\in U$
along any computable unbounded diagonal. First suppose
$\sigma(N)\to\infty$ and set
$x_N=N^{-1}\sum_{n<N}L_n(\sigma(N))$, a computable rational sequence.
The expectation differs from $x_N$ by at most $e_{\sigma(N)}\to0$,
so it cannot converge to $\beta\notin\mathsf D_2$. For samples,
the expected fraction of incorrect registers is at most
$e_{\sigma(N)}$; Markov's inequality gives
$\Pr\{|\rho_N(\sigma(N))-x_N|>a\}\leq e_{\sigma(N)}/a$.
Some deterministic subsequence stays more than $2a$ from $\beta$,
since $x_N$ cannot converge to it. On a sample-convergence event
the displayed deviations exceed $a$ eventually on that subsequence,
an event of probability zero by the bound. A computable unbounded
$\sigma$ has a computably chosen subsequence tending to infinity,
so the same exclusion applies. Bounded-time paths retain the
computable-kernel qualification. An additive density error does not
control logarithms of sparse counts, so the counting-exponent
diagonal statement also retains that qualification.
In the constant-normalizer thermal family of Section~\ref{sec:states},
the same statement is a four-dimensional causal half-space DLR
construction with computable fixed bottom boundary; time is height.
It does not imply unique equilibrium or arbitrary-interaction stability.

\subsection{Null and meager, but dense, exceptional kernels}
\label{subsec:exceptional-kernels}
Here an inquiry is an external computable controller allowed to
inspect successively larger finite spacetime cylinders of the
anchored history, including time-one observations. This access model
is essential for the density conclusion below. It does not assert an
autonomous noisy computer that internally acquires raw-fault data.

\begin{theorem}[Kernel information and fixed-target inquiry]
\label{thm:exceptional-kernels}
Fix an anchored program and a law
$\beta\in\mathsf D_3\setminus\mathsf D_2$ in
Theorem~\ref{thm:robust-third}. At every computable $K\in U$, no
computable inquiry in this model converges to $\beta$ with positive
probability. More generally, every fixed target reached with positive
probability at $K$ is computable in the kernel oracle's jump $K'$.
The kernels in $U$ permitting inquiry to $\beta$ form a null, meager
but dense subset, in intrinsic simplex measure and topology.
\end{theorem}
\begin{proof}
Finite cylinder laws are computable relative to a name of $K$.
Relativize the proof of Theorem~\ref{thm:reach}: on success a target
name is computable in $(X\oplus K)'$. For almost every tape $X$,
relative generalized lowness bounds this by $X\oplus K'$.
The relativized positive-measure cone argument then places the
fixed target below $K'$. This proves the computable-kernel conclusion.

For measure and category, use $q-1$ free probabilities per row;
the last entry is recovered by subtraction. Cover the positive
simplex by rational boxes in these intrinsic coordinates. Rescaling
each box to a cube and interleaving binary digits preserves degrees
off the null, meager dyadic hyperplanes. For almost every code $Y$,
$Y'\equiv_TY\oplus0'$. A positive-measure set on which
$\beta\leq_TK'$ would thus give a positive-measure cone relative
to $0'$, contradicting $\beta\not\leq_T0'$.

For category, 1-generic codes are comeager and also satisfy
$Y'\equiv_TY\oplus0'$. Indeed a generic oracle either meets a
c.e. set of finite halting witnesses or has a prefix with no
extension in it; $Y\oplus0'$ can find which alternative applies.
If a fixed functional computes a non-$Z$-computable set from
$Y\oplus Z$ on a nonmeager Borel set, its success set is comeager
in some cylinder after restriction. Search relative to $Z$ for any
finite extension on which it halts on the requested input. Such an
extension exists, and its open cylinder meets the comeager success
set, so every found answer is correct. This would compute the set
from $Z$. Countably many functionals give the category cone theorem.
Use $Z=0'$ and a cut/name of $\beta$. Thus learnable kernels lie
in a null and meager set, using free coordinates rather than a
fictitious fair-random join of constrained row entries.

For density, in any open patch choose one blank-row probability $p$
in a rational interval $(u,v)$ whose closure stays inside the patch.
Choose rational $u<u_0<v_0<v$ and set
\[
 p=u_0+(v_0-u_0)\sum_{j\geq0}2\,\chi_{0'}(j)4^{-j-1}.
\]
The separated base-four digits give $0'\leq_Tp$. In particular $p$
is noncomputable, hence irrational and non-dyadic. Compensate in a
second entry in the same row; the patch's positive margins allow
this arbitrarily small change. Other free coordinates can be rational.
The sampled entry itself therefore satisfies
$\beta\leq_T0''\leq_Tp'$. At time one, fresh known blank predecessor
neighborhoods away from the finite program give independent samples
of that output, with probability $p$. Let $f_n$ be the frequency
and choose the rational radius $r_n=1/\lceil n^{1/3}\rceil$,
computed by integer comparisons. Hoeffding makes failure of
$[f_n-r_n,f_n+r_n]$ to contain $p$ summable. Eventually these
intervals contain $p$, shrink to it and determine shared binary
prefixes of unbounded length, since $p$ is irrational.
Choose a total $p$-computable rational sequence
$\Theta^p(m)\to\beta$ by the relative limit lemma. Only $p$ is
the oracle: a prefix estimated from one entry is not a name of the
whole kernel. At stage $n$
simulate the first $n$ computations for $n$ steps using the current
finite prefix as a partial oracle, and output the largest completed
index not querying beyond it, or a default. Eventually the prefixes
are correct, and each fixed computation eventually fits and finishes;
the selected indices tend to infinity. Thus this external finite-stage
inquiry converges to $\beta$ almost surely. Every open patch contains
such a kernel.
\end{proof}
The reach set is targetwise existential. This does not prove that
one common learner identifies every unknown law, or that every
$K'$-computable target can be read from every kernel without an
access theorem. Robustness of the law is pointwise on $U$;
unlearnability is not, because the exceptional kernels are dense.

\subsection{A restricted statement about active records}
Strict positive rates give a minimum transition probability
$\varepsilon_{\min}>0$. For any fixed finite set $V$ in a current
layer, every pattern on it has conditional probability at least
$\varepsilon_{\min}^{|V|}$ given the past. Conditional
Borel--Cantelli implies every such pattern recurs infinitely often.
Thus a fixed nonconstant function of that fixed window cannot
converge almost surely. The function and its finite support are fixed,
and it is observed at every integer time. This does not imply that
every convergent numerical estimator has growing support: the one-site
function $X_t(0)/(t+1)$ for a binary spin converges to zero despite
being nonconstant at each time. Constants, expectations and repeated
access to a fixed past record also give no such conclusion. The
separated majorities above have growing support by their definition;
that property is asserted for those readouts only.

